\documentclass[12pt,a4paper]{article}

\usepackage[utf8]{inputenc}
\usepackage[T1]{fontenc}
\usepackage[brazil]{babel}
\usepackage{lmodern}
\usepackage{geometry}
\usepackage{microtype}
\usepackage{array}
\usepackage{tabularx}
\usepackage{enumitem}
\usepackage{xcolor}
\usepackage{fancyhdr}

\geometry{top=1.7cm,bottom=1.7cm,left=1.8cm,right=1.8cm}
\setlength{\parindent}{0pt}
\setlength{\parskip}{0.45em}
\setlist[itemize]{leftmargin=*,itemsep=0.25em,topsep=0.2em}
\definecolor{azul}{HTML}{173B57}
\definecolor{cinza}{HTML}{F2F4F7}

\pagestyle{fancy}
\setlength{\headheight}{15pt}
\fancyhf{}
\lhead{Sequência didática de \textit{machine learning}}
\rhead{Aula 4: k-vizinhos mais próximos}
\cfoot{\thepage}
\renewcommand{\headrulewidth}{0.3pt}

\newcommand{\linha}{\rule{\linewidth}{0.35pt}}
\newcommand{\bloco}[1]{%
  \vspace{0.25em}
  \noindent\colorbox{cinza}{\parbox{0.96\linewidth}{#1}}%
  \vspace{0.3em}
}

\begin{document}

\begin{center}
  {\LARGE\bfseries\color{azul} Aula 4: k-vizinhos mais próximos}\par
  \vspace{0.25em}
  {\large Amostras, proximidade e votação}\par
  \vspace{0.35em}
  {\small Roteiro de aplicação da sequência didática}\par
\end{center}

\bloco{\textbf{Lugar desta aula na pesquisa.} Este é o quarto encontro da sequência de oito aulas. A atividade apresenta uma classificação baseada em exemplos próximos. O registro individual ao final integra o conjunto de evidências previsto para os encontros 2, 4 e 6.}

\section*{1. Identificação}

\begin{tabularx}{\linewidth}{@{}>{\bfseries}p{4.1cm}X@{}}
Turma: & \rule{8.5cm}{0.35pt} \\
Data: & \rule{8.5cm}{0.35pt} \\
Duração: & 50 minutos \\
Organização: & Exploração individual; registro individual ao final. \\
Recursos: & Plataforma no módulo de k-vizinhos mais próximos, computadores do laboratório ou Chromebooks, quadro e folha de registro. \\
\end{tabularx}

\section*{2. Objetivos da aula}

Ao final do encontro, espera-se que os estudantes reconheçam uma classificação baseada em exemplos rotulados, interpretem distância como proximidade no plano e expliquem como a alteração do número de vizinhos pode mudar a votação e a classe atribuída a uma nova amostra.

\textbf{Relação curricular:} EM13CO11 e EM13CO12.

\textbf{Conceitos em foco:} amostra, característica, rótulo, distância, proximidade, vizinho, valor de $k$, votação e classificação.

\section*{3. Preparação antes da aula}

\begin{itemize}
  \item Abrir o módulo de k-vizinhos mais próximos e confirmar o cenário de amostras descritas por doçura e acidez.
  \item Conferir se a nova amostra, as distâncias, os vizinhos selecionados e a votação aparecem de forma legível.
  \item Confirmar que cada estudante terá acesso individual ao módulo, em computador do laboratório ou Chromebook.
  \item Separar uma folha de registro individual para cada estudante e usar o mesmo código anônimo do diagnóstico inicial.
\end{itemize}

\bloco{\textbf{Cuidado didático.} A proximidade no gráfico representa uma relação entre as características escolhidas para a atividade. Não significa que dois alimentos sejam parecidos em todos os aspectos, pois o resultado depende dos dados e da forma de medir distância.}

\section*{4. Desenvolvimento da aula}

\renewcommand{\arraystretch}{1.22}
\begin{tabularx}{\linewidth}{|>{\centering\arraybackslash}p{1.65cm}|p{5.0cm}|X|}
\hline
\textbf{Tempo} & \textbf{Ação com a turma} & \textbf{Mediação docente} \\
\hline
0--7 min & Apresentar uma nova amostra sem rótulo, descrita por doçura e acidez. Retomar a ideia de classificação da aula anterior. & Perguntar quais informações já existem na tabela ou no gráfico e qual informação ainda precisa ser atribuída à nova amostra. \\
\hline
7--13 min & Explorar no quadro a ideia de pontos próximos. Mostrar exemplos já classificados e a nova amostra. & Diferenciar proximidade no plano de uma impressão cotidiana. Cada eixo representa uma característica que foi medida. \\
\hline
13--18 min & Pedir que cada estudante antecipe se a classe atribuída permanece a mesma quando o número de vizinhos aumenta. & Aceitar hipóteses diferentes, desde que indiquem quais exemplos próximos sustentam a previsão. \\
\hline
18--33 min & Explorar o módulo guiado: variar o valor de $k$, observar distâncias, vizinhos selecionados e votos. & Destacar que o modelo consulta exemplos já rotulados. Comparar a votação para dois valores de $k$ sem exigir fórmula de distância. \\
\hline
33--43 min & Retomar previsões individuais e comparar dois valores de $k$ para a mesma amostra. & Solicitar explicação que mencione doçura, acidez, distância e voto, sem depender apenas das cores exibidas no gráfico. \\
\hline
43--50 min & Aplicar e recolher o registro individual. Encerrar com uma síntese oral. & Esclarecer vocabulário sem indicar respostas. Retomar que a classificação depende dos exemplos disponíveis e da regra de vizinhança escolhida. \\
\hline
\end{tabularx}

\section*{5. Perguntas para conduzir a conversa}

\begin{itemize}
  \item ``Que informações descrevem cada amostra?''
  \item ``Qual amostra está mais próxima da nova amostra?''
  \item ``Como os vizinhos participam da decisão?''
  \item ``O que pode mudar quando o valor de $k$ aumenta?''
  \item ``Por que o resultado depende dos exemplos disponíveis?''
\end{itemize}

\section*{6. Evidência a observar}

O registro individual permite verificar se o estudante justifica uma classificação por proximidade e votação, reconhecendo que mudar o valor de $k$ pode mudar os vizinhos considerados e a resposta obtida. No registro docente, anotar explicações que dependam apenas de cores, confusões entre características e rótulos, além de condições de uso da plataforma.

\newpage

\begin{center}
  {\LARGE\bfseries\color{azul} Registro individual}\par
  \vspace{0.2em}
  {\large Aula 4: proximidade, vizinhos e votação}\par
  \vspace{0.35em}
  {\small Material para recolhimento ao final da aula}\par
\end{center}

\bloco{\textbf{Orientação ao estudante.} Esta atividade não vale nota. Responda individualmente e explique como você pensou, mesmo que não tenha certeza. Use o que foi observado na plataforma e o que foi discutido na aula.}

\textbf{Código do participante:} \rule{5cm}{0.4pt}
\hfill
\textbf{Data:} \rule{3cm}{0.4pt}

\section*{Situação}

Um gráfico apresenta alimentos fictícios em dois eixos: doçura e acidez. Alguns exemplos já têm rótulo de categoria A ou categoria B. Uma nova amostra ainda não tem rótulo. O sistema pode observar os exemplos mais próximos para decidir uma categoria.

\section*{Questão 1}

Por que exemplos próximos da nova amostra podem ajudar a escolher uma categoria? Explique usando as características de doçura e acidez.

\vspace{0.25em}
\linha\par\vspace{0.75em}
\linha\par\vspace{0.75em}
\linha

\section*{Questão 2}

Uma nova amostra tem três vizinhos próximos: dois são da categoria A e um é da categoria B. Qual categoria tende a vencer a votação quando $k=3$? Explique.

\vspace{0.25em}
\linha\par\vspace{0.75em}
\linha\par\vspace{0.75em}
\linha

\section*{Questão 3}

Ao aumentar o valor de $k$, outros exemplos mais distantes podem entrar na votação. Como isso pode mudar a categoria atribuída à nova amostra?

\vspace{0.25em}
\linha\par\vspace{0.75em}
\linha

\end{document}
